\documentclass[a4paper,11pt]{article}
\usepackage{exam-style-341}
\examinehead{341 农业综合知识三 · 网络技术与应用}{模拟试卷（二）}
\begin{document}

\examcover{341 农业综合知识三}{网络技术与应用 · 模拟试卷（二）}{网络技术与应用（50 分）}{50 分}{60 分钟}{本科目只考计算机网络，与程序设计、数据库无关。}

\parttitle{网络技术与应用（50 分）}

\qtypename{一、单项选择题}{10 题，每题 2 分，共 20 分}

\q 计算机网络的主要功能是（\quad）
\choices{A. 只能进行数据通信}{B. 数据通信与资源共享}{C. 完全替代邮政通信系统}{D. 只能进行分布式计算}

\q 总线型局域网的拓扑结构中，所有结点都连接到一条公共传输媒体上，其优点是（\quad）
\choices{A. 可靠性高，任一结点故障不影响网络}{B. 结构简单、成本低、易于扩充}{C. 便于故障定位}{D. 便于实现流量控制}

\q 在一个 TCP 连接中，发送方某报文段的序号为 500，共携带 200 字节数据，则接收方正确收到该报文段后返回的确认号应为（\quad）
\choices{A. 500}{B. 700}{C. 701}{D. 200}

\q 关于 IP 地址与 MAC 地址，下列说法正确的是（\quad）
\choices{A. IP 地址随主机在不同网络间移动而改变，MAC 地址不变}{B. MAC 地址随主机在不同网络间移动而改变}{C. 路由器转发分组时改写分组的源 MAC 地址和目的 IP 地址}{D. 路由器的每个接口可以没有 IP 地址}

\q 下列关于分组交换与报文交换的叙述中，正确的是（\quad）
\choices{A. 报文交换不需要存储转发}{B. 分组交换把报文划分成较小的分组，可提高线路利用率和降低时延}{C. 分组交换不产生额外开销}{D. 报文交换一定比分组交换时延小}

\q 在 TCP 的滑动窗口机制中，发送窗口的大小取决于（\quad）
\choices{A. 只取决于接收窗口}{B. 只取决于拥塞窗口}{C. 取接收窗口与拥塞窗口中的较小者}{D. 取接收窗口与拥塞窗口中的较大者}

\q 关于应用层协议所使用的传输层服务，下列说法正确的是（\quad）
\choices{A. FTP 使用 UDP，DNS 只使用 TCP}{B. FTP 使用 TCP，DNS 主要使用 UDP}{C. HTTP 使用 UDP，SMTP 使用 UDP}{D. Telnet 使用 UDP，DNS 使用 TCP}

\q 在 CIDR 中，网络前缀为 \(192.168.16.0/22\) 的地址块共包含多少个可分配给主机的 IP 地址（\quad）
\choices{A. 1024}{B. 1022}{C. 512}{D. 510}

\q 交换机端口收到一个完整帧后先存入缓存，再查表并从相应端口转发出去。下列叙述中正确的是（\quad）
\choices{A. 这个过程会改变帧的内容}{B. 这个过程适用于共享式总线上各站点发送的帧}{C. 转发前交换机可根据帧首部中的目的 MAC 地址查表}{D. 交换机必须把帧广播到所有端口}

\q 在传统以太网中，若争用期为 \(51.2\ \mu s\)、数据率为 10 Mbit/s，则最短帧长为（\quad）
\choices{A. 512 字节}{B. 64 字节}{C. 46 字节}{D. 1500 字节}

\qtypename{二、填空题}{5 题，每题 2 分，共 10 分}

\q 网络协议由语法、语义和\fillblank{}三个要素组成，其中\fillblank{}规定所交换数据的格式。

\q TCP/IP 体系结构从下到上可分为网络接口层、\fillblank{}、运输层和\fillblank{}四层。

\q 在数据链路层的停止等待协议中，发送方每发送一个数据帧就等待对方的\fillblank{}；若经过一段时间仍收不到，就\fillblank{}该数据帧，用序号可以区分新帧与重传帧。

\q 数据传输速率（比特率）是每秒传输的\fillblank{}数；若一个码元携带 \(n\) 个比特，则波特率与比特率的关系是\fillblank{}。

\q TCP 的拥塞控制采用四种算法，即慢开始、拥塞避免、\fillblank{}和\fillblank{}。

\qtypename{三、名词解释}{4 题，每题 2.5 分，共 10 分}

\q 服务访问点 SAP\anslines{1.3cm}

\q MAC 地址\anslines{1.3cm}

\q 网络协议\anslines{1.3cm}

\q 运输层的端口\anslines{1.3cm}

\qtypename{四、简答与计算题}{2 题，每题 5 分，共 10 分}

\q 简述 DNS 域名服务器的分类（按层次划分），并说明主机进行域名解析时本地域名服务器与根域名服务器、权限域名服务器之间的关系。\anslines{4cm}

\q 一个 IP 数据报总长度为 4000 字节（其中首部 20 字节），现要通过 MTU 为 1420 字节的链路传输。请回答：
\begin{enumerate}[leftmargin=2em, itemsep=2pt, topsep=3pt]
  \item 每个分片最多能携带多少字节数据？为什么该值必须是 8 的整数倍？
  \item 该数据报应划分为几个分片？每个分片的片偏移字段取值分别是多少？
  \item 写出各分片的 MF 标志与总长度字段的值（按最小分片数计算）。
\end{enumerate}
\anslines{6cm}

\end{document}
